% ---------------------------------------------------------------------
% Preambolo — tesi giuridica italiana
% ---------------------------------------------------------------------

\usepackage[T1]{fontenc}
\usepackage[utf8]{inputenc}
\usepackage[italian]{babel}
\usepackage{mathptmx}                 % Times: font richiesto dalla gran parte dei regolamenti

\usepackage[a4paper,top=3cm,bottom=3cm,left=3.5cm,right=3cm]{geometry}
\usepackage{setspace}
\onehalfspacing                        % interlinea 1,5

\usepackage{csquotes}
\MakeOuterQuote{"}
\usepackage{emptypage}
\usepackage{titlesec}
\usepackage[bottom,hang,flushmargin]{footmisc}
\usepackage[hidelinks]{hyperref}

% --- Numerazione dei capitoli in numeri romani (Cap. I, II, III...) ---
\renewcommand{\thechapter}{\Roman{chapter}}
\renewcommand{\thesection}{\arabic{section}}
\renewcommand{\thesubsection}{\arabic{section}.\arabic{subsection}}

\titleformat{\chapter}[display]
  {\normalfont\large\bfseries\filcenter}
  {\MakeUppercase{\chaptertitlename}~\thechapter}{1em}{\Large}
\titlespacing*{\chapter}{0pt}{0pt}{40pt}

% --- Note a piè di pagina: numerazione continua in tutta la tesi ------
% In una tesi giuridica la numerazione ricomincia raramente per capitolo.
% Per farla ricominciare, commentare la riga seguente.
\counterwithout{footnote}{chapter}
\renewcommand{\footnotesize}{\small}
\setlength{\footnotesep}{0.9em}

% --- Citazioni testuali lunghe: corpo minore, senza virgolette -------
% Uso: \begin{citazione} ... \end{citazione}
\newenvironment{citazione}
  {\begin{quote}\small\setstretch{1.15}}
  {\end{quote}}

% --- Comandi di servizio ---------------------------------------------
% Marcatori di lavorazione: devono essere ZERO nella versione consegnata.
% Verifica:  grep -n "DAVERIFICARE\|FONTEDAREPERIRE" -r .
\newcommand{\fontedareperire}[1]{\textbf{[FONTE DA REPERIRE: #1]}}
\newcommand{\daverificare}[1]{\textbf{[DA VERIFICARE: #1]}}
\newcommand{\ipotesi}[1]{\textbf{[IPOTESI RICOSTRUTTIVA: #1]}}

% Rinvii interni nella forma giuridica: \supra{sez:cap2:requisiti}
% NOTA: cross_ref_audit.py segnala qui un falso positivo, "\ref{#1} undefined",
% perche' legge il corpo della macro come se fosse un rinvio reale. Va ignorato:
% e' l'unica riscontro atteso su questo template.
\newcommand{\supra}[1]{v.~\emph{supra}, §~\ref{#1}}
\newcommand{\infraref}[1]{v.~\emph{infra}, §~\ref{#1}}

% --- Numerazione delle pagine ----------------------------------------
% La classe report non ha \frontmatter: in main.tex si passa a mano da
% numerazione romana (indice) ad araba (corpo della tesi).
